\documentclass[12pt, a4paper]{article}

% ── Shared preamble (src/lecturenotes.sty) ────────────────────────────────────
%   Packages, colors, boxes, \sol/\soleq and math macros all live in
%   src/*.sty. \input@path points LaTeX two folders up to find them.
\makeatletter
\def\input@path{{../../}}
\makeatother
\usepackage{lecturenotes}

% ── Local macros ──────────────────────────────────────────────────────────────
\providecommand{\bH}{\mathbf{H}}

% ── Local TikZ styles (match S1–S3) ───────────────────────────────────────────
\usetikzlibrary{decorations.pathreplacing}
\tikzset{
  layerlabel/.style={font=\small\itshape, gray},
  opbox/.style={draw=orange!70!black, fill=notebg, thick, rounded corners=4pt,
      minimum height=0.8cm, minimum width=1.7cm, font=\small, align=center},
  convbox/.style={draw=mainblue, fill=boxbg, thick, rounded corners=3pt,
      minimum height=0.75cm, minimum width=1.6cm, font=\small, align=center},
  brace/.style={decorate, decoration={brace, amplitude=4pt}, thick, gray},
}

% \cuboid{<x>}{<depth: channels>}{<height: spatial>}{<color>}{<label>}
\newcommand{\cuboid}[5]{%
  \pgfmathsetmacro{\off}{#3*0.3}
  \filldraw[fill=#4!20, draw=#4!80!black] (#1,-#3/2) rectangle ++(#2,#3);
  \filldraw[fill=#4!35, draw=#4!80!black]
  (#1,#3/2) -- ++(\off,\off) -- ++(#2,0) -- ++(-\off,-\off) -- cycle;
  \filldraw[fill=#4!50, draw=#4!80!black]
  (#1+#2,-#3/2) -- ++(\off,\off) -- ++(0,#3) -- ++(-\off,-\off) -- cycle;
  \node[below, font=\scriptsize, align=center] at (#1+#2/2,-1.75) {#5};
}

% ── Document ──────────────────────────────────────────────────────────────────
\begin{document}

\begin{center}
  {\LARGE \textbf{Section 4.4: Famous CNN Architectures}} \\[0.3em]
  {\large\itshape From LeNet-5 to ResNet, and Transfer Learning} \\[0.5em]
  {\large AI-IM Class --- Lecture Notes}%
  \ifstudentmode
    \quad {\normalsize\color{gray}[Student Copy]}%
  \fi
\end{center}

\vspace{1em}
\hrule
\vspace{1.5em}

With convolution, pooling and the classifier head (Sections 4.2 and 4.3) we
can build any CNN. The question is \emph{which} CNN to build: how many layers,
what kernel sizes, how many channels. The answers came from a short series of
famous architectures, each introducing one idea that is still used today.
In this section we
\begin{itemize}[nosep]
  \item trace LeNet-5, the first successful CNN, layer by layer,
  \item follow the ImageNet competition that made CNNs famous,
  \item study the key idea of AlexNet, VGG, GoogLeNet and ResNet, with a hand
        calculation for each, and
  \item reuse a pretrained network for our own problem (\emph{transfer
          learning}).
\end{itemize}

% ══════════════════════════════════════════════════════════════════════════════
\section{LeNet-5 (1998)}

\textbf{LeNet-5}, by Yann LeCun and colleagues, read handwritten digits on
bank cheques. It already has the structure of Section 4.3: two conv--pool
blocks followed by fully connected layers.

\begin{figure}[H]
  \centering
  \begin{tikzpicture}[>=Stealth, scale=0.95, every node/.style={transform shape}]
    \cuboid{0}{0.1}{3.2}{gray}{input\\$(1,32,32)$}
    \cuboid{1.9}{0.4}{2.8}{mainblue}{C1\\$(6,28,28)$}
    \cuboid{3.9}{0.4}{1.4}{mainblue}{S2\\$(6,14,14)$}
    \cuboid{5.6}{0.9}{1.0}{orange}{C3\\$(16,10,10)$}
    \cuboid{7.9}{0.9}{0.5}{orange}{S4\\$(16,5,5)$}
    \draw[->, thick] (1.1,0) -- (1.8,0);
    \draw[->, thick] (3.1,0) -- (3.8,0);
    \draw[->, thick] (4.8,0) -- (5.5,0);
    \draw[->, thick] (6.9,0) -- (7.8,0);

    \filldraw[fill=green!20, draw=green!50!black] (10.0,-1.2) rectangle ++(0.25,2.4);
    \node[below, font=\scriptsize, align=center] at (10.12,-1.75) {C5\\$120$};
    \filldraw[fill=green!20, draw=green!50!black] (11.3,-0.85) rectangle ++(0.25,1.7);
    \node[below, font=\scriptsize, align=center] at (11.42,-1.75) {F6\\$84$};
    \filldraw[fill=green!35, draw=green!50!black] (12.6,-0.4) rectangle ++(0.25,0.8);
    \node[below, font=\scriptsize, align=center] at (12.72,-1.75) {output\\$10$};
    \draw[->, thick] (9.1,0) -- (9.9,0);
    \draw[->, thick] (10.35,0) -- (11.2,0);
    \draw[->, thick] (11.65,0) -- (12.5,0);
  \end{tikzpicture}
  \caption{LeNet-5. C = convolution ($5\times5$, valid), S = subsampling
    ($2\times2$ average pooling), F = fully connected.}
  \label{fig:lenet}
\end{figure}

The input is a $32 \times 32$ grayscale image (MNIST digits padded from
$28 \times 28$). All convolutions use $5 \times 5$ kernels with stride $1$ and
no padding; all pooling layers are $2 \times 2$ with stride $2$. We use the
common modern version: ReLU activations, average pooling, every C3 filter
connected to all 6 input channels, and a linear output layer.

\begin{stepbox}{Exercise 1 \quad Trace LeNet-5}
  \begin{center}
    \renewcommand{\arraystretch}{1.35}
    \begin{tabular}{llll}
      \toprule
      Layer                 & Operation                    & Output shape                  & Parameters                                     \\
      \midrule
      Input                 &                              & $(1, 32, 32)$                 & ---                                            \\
      C1                    & 6 filters $5\times5$         & $\sol[2.2cm]{(6, 28, 28)}$    & $\sol[4.6cm]{(1\cdot25+1)\cdot6 = 156}$        \\
      S2                    & avg pool $2\times2$          & $\sol[2.2cm]{(6, 14, 14)}$    & $\sol[4.6cm]{0}$                               \\
      C3                    & 16 filters $5\times5$        & $\sol[2.2cm]{(16, 10, 10)}$   & $\sol[4.6cm]{(6\cdot25+1)\cdot16 = 2\,416}$    \\
      S4                    & avg pool $2\times2$          & $\sol[2.2cm]{(16, 5, 5)}$     & $\sol[4.6cm]{0}$                               \\
      C5                    & 120 filters $5\times5$       & $\sol[2.2cm]{(120, 1, 1)}$    & $\sol[4.6cm]{(16\cdot25+1)\cdot120 = 48\,120}$ \\
      F6                    & linear                       & $(84)$                        & $\sol[4.6cm]{(120+1)\cdot84 = 10\,164}$        \\
      Output                & linear                       & $(10)$                        & $\sol[4.6cm]{(84+1)\cdot10 = 850}$             \\
      \midrule
      \multicolumn{3}{l}{Total}                                                            & $\sol[4.6cm]{61\,706}$                         \\
      \bottomrule
    \end{tabular}
  \end{center}
  Why is C5 called a convolution although it behaves like a fully connected
  layer?
  \soleq[1.2cm]{\parbox{0.85\linewidth}{\raggedright Its $5\times5$ kernel
      exactly covers the $5\times5$ input, so there is only one position and
      each filter sees every input value.}}
\end{stepbox}

\begin{notebox}
  \textbf{The original LeNet-5} used scaled $\tanh$ activations, pooling
  layers with one trainable weight and bias per channel, and a hand-designed
  table that connected each C3 filter to only some of the S2 channels
  (1\,516 parameters instead of 2\,416), to save computation in 1998.
  Its total is still about 60 thousand parameters.
\end{notebox}

% ══════════════════════════════════════════════════════════════════════════════
\section{The ImageNet Challenge}

LeNet-5 worked on small grayscale digits. Natural photos are much harder:
thousands of object types, varying lighting, viewpoint and background.
The \textbf{ImageNet Large Scale Visual Recognition Challenge (ILSVRC)} ran
from 2010 to 2017. Its classification task had about 1.2 million training
images in 1\,000 classes. Performance is measured by the \textbf{top-5 error}:
a prediction counts as correct if the true class is among the model's five
most probable classes.

\begin{figure}[H]
  \centering
  \begin{tikzpicture}[x=1.45cm, y=0.13cm]
    \draw[->, gray] (0.3,0) -- (8.7,0);
    \draw[->, gray] (0.3,0) -- (0.3,32) node[above, font=\small] {top-5 error (\%)};
    \foreach \t in {0,10,20,30}
    \node[left, font=\scriptsize] at (0.3,\t) {\t};
    \draw[red!70!black, dashed, thick] (0.4,5.1) -- (8.6,5.1)
    node[right, font=\scriptsize, align=left] {human\\$\approx 5.1$};
    \foreach \i/\yr/\e/\name/\col in {
        1/2010/28.2/{}/gray,
        2/2011/25.8/{}/gray,
        3/2012/16.4/AlexNet/mainblue,
        4/2013/11.7/ZFNet/mainblue,
        5/2014/6.7/GoogLeNet/mainblue,
        6/2015/3.6/ResNet/mainblue,
        7/2016/3.0/{}/mainblue,
        8/2017/2.3/SENet/mainblue} {
        \filldraw[fill=\col!40, draw=\col!80!black] (\i-0.3,0) rectangle (\i+0.3,\e);
        \node[above, font=\scriptsize, fill=white, inner sep=1pt] at (\i,\e+0.4) {\e};
        \node[below, font=\scriptsize] at (\i,0) {\yr};
        \node[below, font=\scriptsize\itshape, gray] at (\i,-2.8) {\name};
      }
  \end{tikzpicture}
  \caption{Winning top-5 error in the ILSVRC classification task. Grey: methods
    based on hand-designed features; blue: CNNs. The human estimate is
    Andrej Karpathy's result on a test sample.}
  \label{fig:ilsvrc}
\end{figure}

The jump in 2012 is the moment deep learning took over computer vision.
By 2015 CNNs beat the human estimate. Each of the architectures below won
(or nearly won) one of these years.

% ══════════════════════════════════════════════════════════════════════════════
\section{AlexNet (2012): Scale It Up}

\textbf{AlexNet} (Krizhevsky, Sutskever and Hinton) is LeNet made bigger:
5 convolutional and 3 fully connected layers, about 60 million parameters,
trained on $224 \times 224$ colour images. Its architecture was not new; what
made it work were practical choices that are still standard today:
\begin{itemize}[nosep]
  \item \textbf{ReLU} instead of $\tanh$, which trained several times faster
        (Section 3.4);
  \item \textbf{GPUs}: two gaming graphics cards, for about a week;
  \item \textbf{data augmentation}: random crops and flips of the training
        images, to get more data for free;
  \item \textbf{dropout} in the fully connected layers, against overfitting
        (we return to it in a later topic).
\end{itemize}

\begin{stepbox}{Exercise 2 \quad AlexNet's first layer}
  The first layer has 96 filters of size $11 \times 11$ with stride $4$ and
  no padding. In the original implementation the input is $(3, 227, 227)$.
  \begin{enumerate}[label=(\alph*), nosep]
    \item Output size: $\sol[5cm]{\lfloor (227 - 11)/4 \rfloor + 1 = 55}$,
          so the output shape is $\sol[3cm]{(96, 55, 55)}$.
    \item Parameters: $\sol[6cm]{(3 \cdot 121 + 1) \cdot 96 = 34\,944}$.
  \end{enumerate}
\end{stepbox}

% ══════════════════════════════════════════════════════════════════════════════
\section{VGG (2014): Only Small Kernels}

\textbf{VGG} (Simonyan and Zisserman, Oxford) asked how far a very simple
design can go. It uses only one kind of convolution, $3 \times 3$ with stride
$1$ and same padding, and only $2 \times 2$ max pooling, repeated in blocks
until the network has 16 or 19 layers. Each block halves the feature maps and
doubles the channels: $64 \to 128 \to 256 \to 512$.

Why only $3 \times 3$? Two stacked $3 \times 3$ layers have the same receptive
field as one $5 \times 5$ layer (Section 4.1), but they are cheaper and have
one more ReLU.

\begin{stepbox}{Exercise 3 \quad Two small kernels vs.\ one large}
  All layers have $C$ input and $C$ output channels. Ignore biases.
  \begin{center}
    \renewcommand{\arraystretch}{1.4}
    \begin{tabular}{lccc}
      \toprule
      Option                     & Receptive field       & Weights            & Weights for $C = 64$ \\
      \midrule
      one $5\times5$ layer       & $5\times5$            & $\sol[1.4cm]{25C^2}$ & $\sol[2cm]{102\,400}$ \\
      two $3\times3$ layers      & $\sol[1.4cm]{5\times5}$ & $\sol[1.4cm]{18C^2}$ & $\sol[2cm]{73\,728}$  \\
      one $7\times7$ layer       & $7\times7$            & $\sol[1.4cm]{49C^2}$ & $\sol[2cm]{200\,704}$ \\
      three $3\times3$ layers    & $\sol[1.4cm]{7\times7}$ & $\sol[1.4cm]{27C^2}$ & $\sol[2cm]{110\,592}$ \\
      \bottomrule
    \end{tabular}
  \end{center}
\end{stepbox}

VGG-16 reached 7.3\% top-5 error, but it has about 138 million parameters,
most of them in its three fully connected layers. Its uniform design made it
popular as a starting point for other tasks.

% ══════════════════════════════════════════════════════════════════════════════
\section{GoogLeNet / Inception (2014): Go Wide, Cheaply}

\textbf{GoogLeNet} (Szegedy et al., Google) won 2014 with 6.7\% top-5 error
using about 7 million parameters, roughly 20 times fewer than VGG. Its
building block is the \textbf{Inception module}: instead of choosing one kernel
size, it applies several in parallel and concatenates the results along the
channel dimension.

\begin{figure}[H]
  \centering
  \begin{tikzpicture}[>=Stealth, x=2.6cm, y=1.15cm]
    \node[opbox, minimum width=3cm] (in) at (1.5, 3) {previous layer};
    \node[convbox] (a1) at (0, 1.5) {$1\times1$};
    \node[convbox] (b1) at (1, 2) {$1\times1$};
    \node[convbox] (b2) at (1, 1) {$3\times3$};
    \node[convbox] (c1) at (2, 2) {$1\times1$};
    \node[convbox] (c2) at (2, 1) {$5\times5$};
    \node[opbox, minimum width=1.6cm] (d1) at (3, 2) {$3\times3$\\[-2pt]\scriptsize max pool};
    \node[convbox] (d2) at (3, 1) {$1\times1$};
    \node[opbox, minimum width=3cm] (out) at (1.5, -0.3) {concatenate channels};
    \draw[->] (in) -- (a1); \draw[->] (in) -- (b1);
    \draw[->] (in) -- (c1); \draw[->] (in) -- (d1);
    \draw[->] (b1) -- (b2); \draw[->] (c1) -- (c2); \draw[->] (d1) -- (d2);
    \draw[->] (a1) -- (out); \draw[->] (b2) -- (out);
    \draw[->] (c2) -- (out); \draw[->] (d2) -- (out);
  \end{tikzpicture}
  \caption{An Inception module. Four branches see the same input with
    different receptive fields. In the second and third branches a
    $1 \times 1$ convolution first shrinks the number of channels, so the
    expensive $3\times3$ and $5\times5$ convolutions become cheap.}
  \label{fig:inception}
\end{figure}

\paragraph{The $1 \times 1$ convolution.}
A $1 \times 1$ kernel looks at a single pixel, but across \emph{all} its
channels. At every pixel it computes a weighted sum of the $C_{\text{in}}$
channel values, which is exactly a small fully connected layer applied at
each pixel separately. It changes the number of channels without touching the
height and width.

\begin{stepbox}{Exercise 4 \quad The $1\times1$ bottleneck}
  \begin{enumerate}[label=(\alph*)]
    \item One pixel has channel values $(2, -1, 3)$. A $1 \times 1$ layer has
          two filters with weights $(1, 0, 1)$ and $(0, 2, -1)$ and no bias.
          The pixel's output is
          $\bigl(\,\sol[1cm]{5},\ \sol[1cm]{-5}\,\bigr)$.
    \item An input of shape $(256, 28, 28)$ goes through a $5 \times 5$
          convolution with 32 filters (same padding). Ignoring biases, it has
          $\sol[5cm]{256 \cdot 25 \cdot 32 = 204\,800}$ weights.
    \item Insert a $1\times1$ convolution with 16 filters first. The weights
          are now
          \[
            \underbrace{\sol[3cm]{256 \cdot 16 = 4\,096}}_{1\times1}
            + \underbrace{\sol[4cm]{16 \cdot 25 \cdot 32 = 12\,800}}_{5\times5}
            = \sol[2cm]{16\,896},
          \]
          about $\sol[1.5cm]{12}$ times fewer. The number of multiplications
          shrinks by the same factor.
  \end{enumerate}
\end{stepbox}

% ══════════════════════════════════════════════════════════════════════════════
\section{ResNet (2015): Skip Connections}

If deeper is better, why not stack 100 layers? He et al.\ (Microsoft
Research) found that plain deep networks did not just overfit: a 56-layer
network had a \emph{higher training error} than a 20-layer one. The deeper
network was harder to \emph{optimise}, even though it could in principle
copy the 20-layer network and set the extra layers to the identity.

Their fix was the \textbf{residual block}. Its layers compute a correction
$F(\bH)$, which is \emph{added} to the block's input through a
\textbf{skip connection} (shortcut):
\begin{resultbox}
  \[
    \text{output} = \relu\bigl(F(\bH) + \bH\bigr).
  \]
  If the best thing a block can do is nothing, it only needs $F = 0$, which
  is easy to learn (all weights near zero).
\end{resultbox}

\begin{figure}[H]
  \centering
  \begin{tikzpicture}[>=Stealth, node distance=0.8cm]
    \node (in) {$\bH$};
    \node[convbox, right=of in] (c1) {$3\times3$ conv};
    \node[opbox, right=0.6cm of c1, minimum width=1.2cm] (r1) {ReLU};
    \node[convbox, right=0.6cm of r1] (c2) {$3\times3$ conv};
    \node[circle, draw, thick, right=1.3cm of c2, inner sep=1pt] (plus) {$+$};
    \node[opbox, right=0.6cm of plus, minimum width=1.2cm] (r2) {ReLU};
    \node[right=0.6cm of r2] (out) {output};
    \draw[->, thick] (in) -- (c1);
    \draw[->, thick] (c1) -- (r1);
    \draw[->, thick] (r1) -- (c2);
    \draw[->, thick] (c2) -- node[below, font=\small] {$F(\bH)$} (plus);
    \draw[->, thick] (plus) -- (r2);
    \draw[->, thick] (r2) -- (out);
    \draw[->, thick, red!70!black] ($(in.east)+(0.2,0)$) -- ++(0,-1.1)
    -| node[pos=0.25, below, font=\small] {identity (skip connection)} (plus);
    \draw[brace] ($(c1.north west)+(0,0.15)$) --
    node[above=5pt, font=\small] {$F$: the residual branch}
    ($(c2.north east)+(0,0.15)$);
  \end{tikzpicture}
  \caption{A residual block. The input skips past two convolutions and is
    added to their output. (Real ResNets also put a normalisation layer after
    each convolution, which we meet in a later topic.)}
  \label{fig:resblock}
\end{figure}

\paragraph{Why it helps the gradient.}
Recall the vanishing-gradient problem of Section 3.4: backpropagation
multiplies one derivative per layer, and a product of many small numbers is
tiny. Take one scalar ``layer'' $y = f(x)$ with derivative $f'$. The residual
version $y = x + f(x)$ has derivative $1 + f'$. The $1$ comes from the
skip connection, and it does not shrink.

\begin{stepbox}{Exercise 5 \quad Gradient through three layers}
  Each layer has $f' = 0.1$. Compute $\partial y / \partial x$ through three
  layers in a row.
  \[
    \text{plain: } (f')^3 = \sol[2cm]{0.001},
    \qquad
    \text{residual: } (1 + f')^3 = \sol[2cm]{1.331}.
  \]
  Through 50 plain layers the gradient would be $0.1^{50} = 10^{-50}$.
\end{stepbox}

With skip connections, He et al.\ trained networks with 152 layers and won
2015 with 3.6\% top-5 error, better than the human estimate. When the two
shapes differ (the block changes the number of channels or uses stride $2$),
the shortcut uses a $1 \times 1$ convolution to match them. Skip connections
are now everywhere in deep learning, including the transformers behind
large language models.

% ══════════════════════════════════════════════════════════════════════════════
\section{Comparison}

\begin{center}
  \renewcommand{\arraystretch}{1.35}
  \begin{tabular}{lcrrrl}
    \toprule
    Network      & Year & Layers & Parameters & Top-5 error & Key idea                      \\
    \midrule
    LeNet-5      & 1998 & 5      & 60\,K      & ---         & conv + pool + dense           \\
    AlexNet      & 2012 & 8      & 60\,M      & 16.4\%      & ReLU, GPUs, augmentation      \\
    VGG-16       & 2014 & 16     & 138\,M     & 7.3\%       & stacks of $3\times3$ kernels  \\
    GoogLeNet    & 2014 & 22     & 7\,M       & 6.7\%       & Inception, $1\times1$ bottleneck \\
    ResNet-152   & 2015 & 152    & 60\,M      & 3.6\%       & skip connections              \\
    \bottomrule
  \end{tabular}
\end{center}

\begin{notebox}
  \textbf{After 2015.} Later designs mostly refine these ideas.
  \emph{MobileNet} splits each convolution into a per-channel spatial step and
  a $1\times1$ step so that it runs on phones and embedded cameras.
  \emph{EfficientNet} scales depth, width and image resolution together.
  \emph{Vision Transformers} (2020) replace convolution with attention, and
  \emph{ConvNeXt} (2022) showed that a modernised ResNet is just as good.
  CNNs remain the default choice when data or computing power is limited,
  as in most factory applications.
\end{notebox}

% ══════════════════════════════════════════════════════════════════════════════
\section{Transfer Learning}

Training a ResNet from scratch needs a million labelled images and days of
GPU time. A factory that wants to classify, say, three kinds of surface defect
may have only a few hundred photos. \textbf{Transfer learning} solves this by
reusing a network that was already trained on ImageNet.

The idea follows from Section 4.3. The feature extractor of an ImageNet
network has learned general features (edges, textures, shapes) that are
useful for almost any image. Only the classifier head is specific to the 1\,000
ImageNet classes. So we
\begin{enumerate}[nosep]
  \item load the pretrained network,
  \item replace its final linear layer with a new one that has one output per
        \emph{our} class, and
  \item train on our data: either only the new layer (\emph{feature
          extraction}, with the rest frozen), or the whole network with a small
        learning rate (\emph{fine-tuning}).
\end{enumerate}

\begin{codebox}
import torchvision

model = torchvision.models.resnet18(weights="IMAGENET1K_V1")
for p in model.parameters():
    p.requires_grad = False              # freeze the feature extractor

model.fc = nn.Linear(512, 3)             # new head: 3 defect classes
optimizer = torch.optim.Adam(model.fc.parameters(), lr=1e-3)
\end{codebox}

\begin{stepbox}{Exercise 6 \quad How much is trained?}
  ResNet-18 has about 11.7 million parameters. Its final global average pooling
  produces 512 features.
  \begin{enumerate}[label=(\alph*), nosep]
    \item The new head has $\sol[4cm]{(512 + 1) \cdot 3 = 1\,539}$
          parameters.
    \item With the feature extractor frozen, we train about
          $\sol[3cm]{0.013\%}$ of the network.
  \end{enumerate}
\end{stepbox}

% ══════════════════════════════════════════════════════════════════════════════
\section{Summary}

\begin{itemize}
  \item LeNet-5 set the pattern: conv--pool blocks followed by dense layers.
        Tracing its shapes and parameters uses only the formulas of
        Sections 4.2 and 4.3.
  \item AlexNet (2012) showed that a large CNN trained on GPUs with ReLU,
        augmentation and dropout beats hand-designed features by a wide
        margin.
  \item VGG: stacks of $3\times3$ kernels give large receptive fields with
        fewer weights and more nonlinearity.
  \item GoogLeNet: parallel branches, with $1\times1$ convolutions to cut the
        number of channels cheaply.
  \item ResNet: skip connections, $\text{output} = \relu(F(\bH) + \bH)$, keep
        gradients alive and make networks of 100+ layers trainable.
  \item Transfer learning reuses an ImageNet feature extractor and retrains
        only a new head, so a few hundred images can be enough.
\end{itemize}

\paragraph{Coming up.}
Section 4.5 goes beyond classifying a whole image: finding \emph{where}
objects are (detection, segmentation) and reading text (OCR).

\end{document}
